\documentclass[10pt,a4paper]{article}

\usepackage[utf8]{inputenc}
\usepackage[T1]{fontenc}
\usepackage{times}
\usepackage{color}
\usepackage{graphicx}
\usepackage{url}
\usepackage[hidelinks,colorlinks=true,linkcolor=navy,citecolor=navy,urlcolor=navy]{hyperref}

% Math text compatibility macro for vanilla LaTeX
\providecommand{\text}[1]{\mbox{#1}}
\providecommand{\mathrm}[1]{\textrm{#1}}

% Page Geometry (A4 margins 1.5 cm)
\setlength{\topmargin}{-1.5cm}
\setlength{\headheight}{12pt}
\setlength{\headsep}{12pt}
\setlength{\textheight}{26.2cm}
\setlength{\footskip}{18pt}
\setlength{\oddsidemargin}{-0.8cm}
\setlength{\evensidemargin}{-0.8cm}
\setlength{\textwidth}{18.0cm}
\setlength{\parindent}{0pt}
\setlength{\parskip}{2.5pt plus 0.8pt minus 0.5pt}

% Color Definitions
\definecolor{navy}{rgb}{0.05, 0.20, 0.55}
\definecolor{teal}{rgb}{0.05, 0.45, 0.45}
\definecolor{darkred}{rgb}{0.65, 0.10, 0.10}
\definecolor{grayborder}{rgb}{0.75, 0.75, 0.75}
\definecolor{boxbg}{rgb}{0.95, 0.97, 1.00}
\definecolor{alertbg}{rgb}{0.94, 0.98, 0.95}

% Custom Box Command
\newcommand{\academicbox}[2]{%
  \begin{center}%
  \setlength{\fboxsep}{5pt}%
  \setlength{\fboxrule}{0.8pt}%
  \fcolorbox{navy}{boxbg}{%
    \parbox{17.4cm}{%
      \textbf{\color{navy}#1}\\[2pt]%
      #2%
    }%
  }%
  \end{center}%
}

\newcommand{\alertbox}[2]{%
  \begin{center}%
  \setlength{\fboxsep}{5pt}%
  \setlength{\fboxrule}{0.8pt}%
  \fcolorbox{teal}{alertbg}{%
    \parbox{17.4cm}{%
      \textbf{\color{teal}#1}\\[2pt]%
      #2%
    }%
  }%
  \end{center}%
}

\begin{document}

% Header Block
\begin{center}
    {\Large \textbf{\color{navy}LAPORAN KAJIAN AKADEMIK KOMPREHENSIF DAN METODOLOGI}}\\[2pt]
    {\large \textbf{\color{navy}Sistem Deteksi Aktivitas Lapar Bibit Lele (\textit{Clarias sp.}) Berbasis Visi Komputer Edge}}\\[2pt]
    {\normalsize \textbf{Pemisahan Riak Ikan vs. Derau Aerator Kontinu, Reduksi Gangguan Luar Ruangan, Formulasi $FAI$, \& Kendali Servo}}\\[4pt]
    \rule{\textwidth}{1.0pt}\\[2pt]
    \footnotesize
    \begin{tabular}{p{3.2cm} p{5.5cm} p{3.2cm} p{5.5cm}}
        \textbf{Nama Mahasiswa} & : Khoirotun Nisa' & \textbf{Dosen Pembimbing} & : Dr. Ulla Delfana Rosiani, ST., MT. \\
        \textbf{NIM} & : 2341720057 & \textbf{Program Studi} & : Sarjana Terapan Teknik Informatika \\
        \textbf{Institusi} & : Politeknik Negeri Malang & \textbf{Objek Riset} & : Bibit Ikan Lele Fase Pendederan ($2\text{--}7\text{ cm}$) \\
        \textbf{Kondisi Kolam} & : Terpal Bulat Outdoor + Aerator & \textbf{Arsitektur Edge} & : Raspberry Pi 5 + Logitech Brio 100 + PCA9685 \\
    \end{tabular}\\[1pt]
    \rule{\textwidth}{1.0pt}
\end{center}

\vspace{-0.3cm}

\academicbox{RINGKASAN EKSEKUTIF \& KECUKUPAN LITERATUR ILMIAH}{
Dokumen ini disusun sebagai panduan akademik dan teknis terlengkap untuk dipelajari, dipahami, dan dicetak sebagai rujukan utama penyusunan proposal dan naskah skripsi Bab 2 (Tinjauan Pustaka) dan Bab 3 (Metodologi Penelitian).
\begin{itemize}
    \item \textbf{Status Kecukupan Literatur (Sangat Cukup \& Solid)}: Sebanyak \textbf{20 artikel Computer Vision} internasional (termasuk jurnal bereputasi 2024--2026) dan \textbf{14 artikel nutrisi/etologi bibit lele} yang telah dihimpun sudah \textbf{sangat memadai dan tuntas}. Peneliti \textbf{TIDAK PERLU mencari literatur baru via Consensus} agar tidak terjadi pemborosan fokus (*information overload*).
    \item \textbf{Kunci Implementasi Nyata}: Fokus skripsi saat ini adalah: (1) pemodelan baseline aerator kontinu tanpa mematikan aerasi, (2) eliminasi derau dinding terpal biru dan pantulan cahaya matahari, (3) kuantisasi skor $FAI$, dan (4) integrasi kuota pakan harian pada servo PCA9685.
\end{itemize}
}

\vspace{-0.2cm}

\section*{1. Analisis Visual Setup Aktual Kamera di Atas Permukaan Air}

\begin{minipage}[t]{0.35\textwidth}
    \vspace{0pt}
    \centering
    \includegraphics[width=\textwidth]{figures/kolam-setup.png}\\[3pt]
    \parbox{\textwidth}{\centering \scriptsize \textbf{Gambar 1}: Setup Kamera Jarak Dekat ($h \approx 30\text{--}50\text{ cm}$) di Atas Kolam Terpal Ber-Aerator.}
\end{minipage}
\hfill
\begin{minipage}[t]{0.62\textwidth}
    \vspace{0pt}
    \textbf{1.1 Karakteristik Optik Berjarak Sangat Dekat (\textit{Close-Range Nadir})}\\[2pt]
    \small Sebagian besar literatur menempatkan kamera pada tiang tinggi ($h = 1,2\text{--}3,0\text{ m}$, seperti Qian et al., 2025; Hu et al., 2026), di mana riak halus benih lele melebur menjadi blur. Pada setup mahasiswa:
    \begin{enumerate}\setlength{\itemsep}{1pt}\setlength{\parskip}{0pt}
        \item \textbf{Ketinggian Rendah ($h \approx 30\text{--}50\text{ cm}$)}: Resolusi spasial tinggi; buih aerator ($2\text{--}10\text{ mm}$) dan sambaran benih lele kecil ($2\text{--}7\text{ cm}$) terekam tajam puluhan piksel.
        \item \textbf{Sudut Nadir Tegak Lurus ($\theta \approx 80^\circ\text{--}90^\circ$)}: Mengeliminasi distorsi perspektif miring (\textit{keystone distortion}).
        \item \textbf{Pemotongan Bidang Minat (\textit{Static ROI Masking})}: Dinding terpal biru pada $30\%$ bagian atas frame wajib dibuang melalui mask statis agar getaran terpal akibat hembusan angin luar ruangan tidak memicu deteksi palsu.
    \end{enumerate}
\end{minipage}

\vspace{0.15cm}

\textbf{1.2 Dekomposisi Fenomena Visual Air Keruh \& Buih Aerator Cabang Y}
\begin{itemize}\setlength{\itemsep}{0.5pt}\setlength{\parskip}{0pt}
    \item \textbf{Kekeruhan Pekat (*High Turbidity*)}: Air kolam pendederan lele berwarna hijau kecokelatan pekat dengan kedalaman tembus pandang $\le 3\text{ cm}$. Pendekatan deteksi bentuk fisik tubuh ikan (*fish bounding box*) dipastikan \textbf{gagal total}. Deteksi wajib berfokus pada \textbf{dinamika riak gelombang dan semburan buih sesaat di permukaan air}.
    \item \textbf{Semburan Gelembung Aerator Kontinu}: Percabangan selang Y melepaskan gelembung udara terus-menerus. Gelembung membentuk bintik putih berkilau (*specular highlights*) yang melayang keluar dan pecah. Gejolak ini merupakan \textbf{baseline pergerakan normal kolam}.
\end{itemize}

\newpage

\section*{2. Sintesis Kritis 20 Artikel Literatur Global \& Penetapan \textit{Research Gap}}

Tabel 1 menyajikan ekstraksi komparatif terhadap 20 literatur ilmiah visi komputer akuakultur internasional yang telah dihimpun di repositori riset.

\subsection*{Tabel 1: Matriks Komparasi 20 Literatur Terhadap Parameter Skripsi Mahasiswa}
\vspace{-0.2cm}
{\scriptsize
\setlength{\tabcolsep}{2pt}
\begin{center}
\begin{tabular}{|p{0.35cm}|p{2.2cm}|p{2.0cm}|p{1.5cm}|p{2.3cm}|p{2.0cm}|p{1.8cm}|p{4.7cm}|}
\hline
\textbf{No} & \textbf{Penulis \& Thn} & \textbf{Objek \& Air} & \textbf{Setup \& Posisi} & \textbf{Metode \& Model} & \textbf{Penanganan Noise} & \textbf{Hasil Evaluasi} & \textbf{Keterbatasan Kritis Terhadap Skripsi Mahasiswa} \\
\hline
01 & Qian et al. (2025) & Ikan mas, outdoor & Top, $h=1,8\text{ m}$ & PM-YOLO (YOLOv8n) & PConv clutter filter & Prec 87.7\%, mAP 83.5\% & Mengandalkan kontras fisik tubuh di air jernih; gagal di kolam lele keruh pekat. \\
\hline
02 & Hu et al. (2026) & Bass, RAS tank & Miring, $h=0,7\text{ m}$ & HVIT (ResNet+ViT+LSTM) & Self-attention ViT & Akurasi 95.10\%, 145 ms & Model raksasa (28,6M param), butuh GPU RTX 3090; tidak bisa jalan di CPU Raspi. \\
\hline
03 & Alahmad et al. (2025) & Seabream, laut & Submersible $3,4\text{ m}$ & YOLOv8 + Optical Flow & Velocity thresholding & Akurasi deteksi 85\% & Menuntut kamera underwater kedap air \& laut jernih; tidak cocok kolam terpal. \\
\hline
04 & Fukae et al. (2023) & Ikan air tawar & Top, $h=1,5\text{ m}$ & Farnebäck Optical Flow & Magnitude threshold & Akurasi 88.3\% & Farnebäck dense flow sangat lambat di Raspi (<2 FPS); gagal jika aerator aktif kontinu. \\
\hline
05 & Zhao et al. (2025) & Koki, tangki jernih & Top, $h=1,2\text{ m}$ & Keypoint Pose + ST-GCN & Topologi graph anatomi & Akurasi feeding 98.5\% & Menuntut air bening agar sendi ikan terbaca; gagal pada air pendederan keruh. \\
\hline
06 & Li et al. (2020) & Multi-spesies & Tinjauan multi & Systematic CV Review & Filter polarisasi optik & Sintesis komparatif & Menegaskan bahwa pemodelan noise buih aerator di kolam outdoor belum terpecahkan. \\
\hline
07 & Marliza \& Aisuwarya (2021) & Lele, kolam terpal & Top, $h\approx 1,0\text{ m}$ & Gray + Binary Threshold & Tanpa filter aerator & Pakan saat kumpul & Sangat rentan \textit{false alarm} oleh aerator dan terik matahari; tanpa kuota pakan harian. \\
\hline
08 & Atoum et al. (2015) & Kawanan, indoor & Top, $h=1,0\text{ m}$ & Dynamic Texture Energy & Normalisasi kontras & Feeding stops at rest & Mengasumsikan kolam tenang tanpa aerator aktif terus-menerus di area pandang kamera. \\
\hline
09 & Al-Abri et al. (2025) & Review akuakultur & Review global & YOLOv7--v9, Transformer & Domain adaptation & Sintesis SOTA & Menyoroti disparitas ekstrem antara kompleksitas model DL dan komputasi edge. \\
\hline
10 & Zhou et al. (2018) & Nila, indoor & Top, $h=1,2\text{ m}$ & Near-Infrared + ANFIS & Filter optik NIR & Akurasi feeding 98.0\% & Membutuhkan proyektor inframerah khusus; tidak efektif di bawah sinar terik matahari. \\
\hline
11 & Liu et al. (2014) & Salmon, indoor & Top, $h=1,5\text{ m}$ & Frame Difference (CVFAI) & Ambang piksel minimal & Korelasi nafsu $R^2 > 0.90$ & Asumsi dasar air tenang; gelembung aerator akan selalu salah terdeteksi sebagai pakan. \\
\hline
12 & Zhao et al. (2017) & Nila, indoor & Top, $h=1,2\text{ m}$ & Trajectory + Kinetic Energy& Tracking filter Kalman & Akurasi feeding 92.5\% & Menuntut pelacakan individu di air bening; bibit lele berenang berkerumun pekat. \\
\hline
13 & Zhao et al. (2016) & Ikan mas, terkontrol & Top, $h=1,2\text{ m}$ & BMFS + Partition Flow & Adaptive YCbCr & Korelasi agresi tinggi & Optical flow per-piksel sangat berat di edge; tidak memodelkan letupan aerator periodik. \\
\hline
14 & Dong et al. (2024) & Kerapu, terkontrol & Top, $h=0,5\text{ m}$ & Three-Stream CNN & Multi-stream fusion & Akurasi 99.28\% & Tiga cabang CNN simultan mustahil dieksekusi secara real-time pada CPU Raspberry Pi. \\
\hline
15 & Zhang et al. (2024) & Bawah laut dinamis & Submersible horisontal & RUSNet Segmentation & Temporal blur selection & mIoU 92.4\% video buram & Khusus segmentasi bawah air laut; tidak relevan untuk kuantifikasi riak pakan overhead. \\
\hline
16 & Aris et al. (2023) & Gerak umum & CCTV sudut miring & Gaussian Mixture Model & Multi-Gaussian latar & Subtraction evaluation & GMM standar gagal pada aerator kontinu (terjadi \textit{cavity} \& false foreground berkepanjangan). \\
\hline
17 & Kamal et al. (2025) & Lele ber-aerator & Hidrofon akustik & Audio ML (SVM, RF) & Bandpass filter suara & Akurasi kunyah 91.2\% & Bising getaran kompresor aerator merusak SNR audio; metode visi komputer lebih handal. \\
\hline
18 & Yang et al. (2025) & Gelombang pakan & Review multi-sudut & Spatio-Temporal Transformer & Attention channel & Analisis tren 2025 & Menegaskan krusialnya pemodelan deret waktu (*temporal*) untuk memisahkan gelombang air. \\
\hline
19 & Vijayalakshmi (2025) & Ikan air tawar out. & Top \& Oblique & AquaYOLO (Enhanced v8) & CSP layer feature & mAP 90.9\% & Bounding box gagal saat bibit lele bersembunyi di bawah buih aerator atau air keruh. \\
\hline
20 & Wu et al. (2024) & Bass, indoor & Top, $h=1,2\text{ m}$ & DeformAtt-ViT (Splash) & Splash thumbnail crop & Akurasi 95.50\%, F1 94.13\% & \textbf{Gagal pada bibit lele}; cipratan benih halus tidak menghasilkan buih putih masif bass dewasa. \\
\hline
-- & \textbf{Skripsi Nisa'} & \textbf{Bibit Lele Outdoor} & \textbf{Top Dekat, $h \approx 40\text{ cm}$} & \textbf{Adaptive Baseline + FAI} & \textbf{$Z$-Score + HSV Filter} & \textbf{Target: Akurasi $\ge 85\%$} & \textbf{SOLUSI TEPAT: Ringan di Raspi, atasi aerator kontinu, air keruh, \& kuota terikat.} \\
\hline
\end{tabular}
\end{center}
}

\newpage

\subsection*{2.2 Formulasi Empat Kesenjangan Ilmiah Berbasis Bukti (\textit{Research Gaps})}
\begin{enumerate}
    \item \textbf{Gap 1: Kegagalan Pengenalan Morfologi Fisik Tubuh pada Air Pendederan Keruh}:
    Riset modern (Qian et al., 2025; Zhao et al., 2025; Vijayalakshmi, 2025) mengandalkan deteksi siluet tubuh ikan. Pada bibit lele di air keruh, tubuh ikan tidak tampak kecuali saat moncongnya memecah permukaan air. \textbf{Kebaruan}: Menggeser paradigma dari deteksi objek tubuh menjadi \textbf{analisis dinamika fluktuasi riak gelombang permukaan air}.
    \item \textbf{Gap 2: Asumsi Naif Permukaan Air Tenang Tanpa Gejolak Aerator}:
    Metode klasik (Liu et al., 2014; Marliza \& Aisuwarya, 2021; Atoum et al., 2015) menganggap kolam berada pada kondisi tenang tanpa gerak sebelum pakan diberikan. Pada kolam pendederan lele nyata, aerator wajib menyala terus-menerus untuk menjaga pasokan oksigen terlarut ($DO > 3\text{ mg/L}$). Tanpa pemodelan baseline, sistem konvensional akan selalu memicu pakan (*false alarm 100\%*). \textbf{Kebaruan}: Memodelkan derau aerator sebagai \textbf{baseline dinamis berbasis statistik berjalan (*running Gaussian*)}, sehingga pakan hanya terpicu jika ada deviasi positif di atas batas varians aerator.
    \item \textbf{Gap 3: Ketidaksesuaian Asumsi Cipratan Kasar Ikan Dewasa vs. Riak Halus Bibit}:
    Penelitian Wu et al. (2024) mengandalkan buih putih tebal kibasan ekor bass dewasa (*splash thumbnails*). Benih lele ($2\text{--}7\text{ cm}$) tidak menghasilkan busa putih tebal. \textbf{Kebaruan}: Menempatkan kamera sangat dekat ($30\text{--}50\text{ cm}$) dan memanfaatkan \textbf{Three-Frame Differencing terbobot} untuk menangkap riak mikro halus benih secara presisi tinggi.
    \item \textbf{Gap 4: Disparitas Kebutuhan Komputasi Deep Learning vs. Daya Edge Device}:
    Model mutakhir Hu et al. (2026) dan Dong et al. (2024) menuntut workstation GPU puluhan juta rupiah. \textbf{Kebaruan}: Merancang arsitektur komputasi matematika deterministik yang mampu dieksekusi secara efisien pada CPU Raspberry Pi 5 pada kecepatan stabil 10 FPS.
\end{enumerate}

\section*{3. Metodologi Membedakan Gerakan Bibit Lele vs. Derau Buih Aerator}

\subsection*{3.1 Karakterisasi Spasial dan Temporal Riak Air}
Tabel 2 memetakan perbedaan tanda fisik (*physical signatures*) antara letupan buih aerator dan sambaran bibit lele.

\begin{table}[h!]
\centering
\small
\setlength{\tabcolsep}{3pt}
\caption{Matriks Perbedaan Karakteristik Fisik: Buih Aerator vs. Riak Makan Bibit Lele}
\label{tab:karakteristik-gerak}
\begin{tabular}{|p{3.2cm}|p{6.8cm}|p{6.8cm}|}
\hline
\textbf{Parameter Fisik} & \textbf{Derau Buih Aerator (Baseline Normal)} & \textbf{Riak Sambaran Bibit Lele (Sinyal Lapar)} \\
\hline
\textbf{Distribusi Spasial} & Terpusat pada titik keluaran selang Y (*point source*), melayang keluar secara laminar. & Tersebar acak di seluruh permukaan kolam (*stochastic multi-point distribution*). \\
\hline
\textbf{Karakteristik Temporal} & Kontinu, periodik, osilasi frekuensi relatif konstan ($f \approx 1\text{--}2\text{ Hz}$). & Transien, mendadak (*burst*), berdurasi singkat 3--6 frame ($0,3\text{--}0,6\text{ detik}$ pada 10 FPS). \\
\hline
\textbf{Stabilitas Varians Piksel} & Varians intensitas pergerakan rendah dan stabil sepanjang waktu. & Lonjakan tajam energi kinetik yang melampaui rentang varians normal ($Z > 3,5$). \\
\hline
\textbf{Morfologi Optik} & Gelembung spekular berbentuk cakram bulat kecil yang meletup ke luar. & Gelombang cincin riak air mikro konsentris yang saling berinterferensi. \\
\hline
\end{tabular}
\end{table}

\subsection*{3.2 Algoritma Pemodelan Baseline Berjalan (\textit{Welford Running Gaussian Baseline})}
Untuk mengisolasi sinyal tanpa mematikan aerator, sistem memperbarui nilai rata-rata ($\mu_{\text{base}}$) dan varians ($\sigma_{\text{base}}^2$) dari pergerakan kolam secara rekursif:
\begin{equation}
\mu_{\text{base}}(t) = \alpha \cdot M_t + (1 - \alpha) \cdot \mu_{\text{base}}(t-1)
\end{equation}
\begin{equation}
\sigma_{\text{base}}^2(t) = \alpha \cdot (M_t - \mu_{\text{base}}(t))^2 + (1 - \alpha) \cdot \sigma_{\text{base}}^2(t-1)
\end{equation}
Di mana $M_t$ adalah persentase piksel aktif pada frame ke-$t$, dan $\alpha$ adalah laju adaptasi (*learning rate*, $\alpha \approx 0,02\text{--}0,05$) agar lonjakan sesaat akibat pakan tidak merusak model baseline normal aerator.

\subsection*{3.3 Uji Deviasi Baku ($Z$-Score Anomaly Detection) \& Verifikasi Temporal}
Setiap frame yang masuk dievaluasi tingkat deviasinya terhadap model baseline:
\begin{equation}
Z_t = \frac{M_t - \mu_{\text{base}}(t)}{\sigma_{\text{base}}(t)}
\end{equation}
Kriteria keputusan ditetapkan sebagai berikut:
\begin{itemize}
    \item \textbf{Kondisi Normal Aerator ($Z_t \le 2,0$)}: Fluktuasi gerak berada dalam batas 2 standar deviasi letupan aerator normal $\to$ \textbf{Diabaikan sebagai Noise}.
    \item \textbf{Kondisi Sambaran Lapar ($Z_t \ge 3,5$)}: Terjadi anomali pergerakan air signifikan $\to$ \textbf{Kandidat Sinyal Pakan}.
    \item \textbf{Verifikasi Durasi Temporal ($N \ge 3$ Frame)}: Untuk mencegah pemicuan akibat satu gelembung besar aerator yang meletup keras pada satu frame tertentu, sistem mensyaratkan kondisi $Z_t \ge 3,5$ harus bertahan selama \textbf{minimal 3 hingga 5 frame berturut-turut ($0,3\text{--}0,5\text{ detik}$)}.
\end{itemize}

\newpage

\section*{4. Metodologi Reduksi Derau Lingkungan Luar Ruangan (\textit{Outdoor Noise})}

\subsection*{4.1 Eliminasi Pantulan Kilau Matahari (\textit{Sun Glint \& Specular Reflection})}
\begin{enumerate}
    \item \textbf{Transformasi Ruang Warna HSV}: Frame RGB dikonversikan ke ruang warna HSV ($\mathbf{I}_{\text{HSV}} = \{H, S, V\}$). Komponen kecerahan $V$ yang rentan terhadap perubahan cahaya dinormalkan menggunakan CLAHE (*Contrast Limited Adaptive Histogram Equalization*). Operasi selisih frame diprioritaskan pada kanal saturasi $S$, karena pantulan sinar matahari selalu mengalami desaturasi mendekati putih netral ($S \approx 0$).
    \item \textbf{Pemasangan Filter Polarisasi Sirkular (CPL)}: Sesuai rujukan Li et al. (2020), filter optik polarisasi (CPL) dipasang di depan lensa kamera Logitech Brio 100. Filter ini secara optik memblokir gelombang pantulan air hingga $80\%$, menghilangkan kilau matahari tanpa membebani komputasi CPU Raspberry Pi.
\end{enumerate}

\subsection*{4.2 Eliminasi Daun Rontok \& Serangga Mengapung}
Daun rontok bergerak pasif secara lambat ($v < 5\text{ cm/s}$) dan satu arah (*unidirectional drift*) mengikuti hembusan angin atau arus air. Sebaliknya, riak sambaran bibit lele bergerak bolak-balik frekuensi tinggi secara multi-arah (*omnidirectional*). Vektor gerak lambat satu arah difilter keluar.

\subsection*{4.3 Eliminasi Bayangan Burung Melintas atau Petani}
Bayangan eksternal membentuk satu gumpalan kontur raksasa (*large silhouette blob*) yang menutupi $> 20\%$ area permukaan kolam secara serentak. Sebaliknya, riak bibit lele membentuk klaster riak mikro ($50\text{--}500\text{ piksel}$). Setiap kontur gerak dengan luas $\text{Area} > 0,15 \times \text{Area}_{\text{ROI}}$ otomatis diabaikan sebagai bayangan luar.

\section*{5. Kuantifikasi Perubahan Piksel Menjadi Feeding Activity Index ($FAI$)}

\subsection*{5.1 Three-Frame Differencing (Menghilangkan Efek \textit{Ghosting})}
Pengurangan dua frame citra standar ($|I_t - I_{t-1}|$) rentan menghasilkan bayangan semu (*ghosting*). Untuk menghasilkan kontur tepi riak yang presisi, diterapkan selisih tiga frame berurutan:
\begin{equation}
\Delta_1(x,y) = |I_t(x,y) - I_{t-1}(x,y)|, \quad \Delta_2(x,y) = |I_{t-1}(x,y) - I_{t-2}(x,y)|
\end{equation}
\begin{equation}
D_{\text{motion}}(x,y) = \mathcal{T}_{\text{Otsu}}(\Delta_1(x,y)) \cap \mathcal{T}_{\text{Otsu}}(\Delta_2(x,y))
\end{equation}
Operasi logika biner \textit{bitwise AND} ($\cap$) memastikan bahwa piksel aktif yang dipertahankan hanyalah piksel yang benar-benar mengalami pergerakan pada posisi seketika frame ke-$t$.

\subsection*{5.2 Kalkulasi Energi Riak ($M_t$) dan Skor $FAI_t$}
Total persentase piksel aktif dihitung di dalam wilayah minat air ($\text{ROI}_{\text{air}}$):
\begin{equation}
M_t = \frac{\sum_{(x,y) \in \text{ROI}_{\text{air}}} D_{\text{motion}}(x,y)}{\text{Jumlah Total Piksel}(\text{ROI}_{\text{air}})} \times 100\%
\end{equation}
\begin{equation}
FAI_t = \frac{M_t - \mu_{\text{base}}(t)}{\sigma_{\text{base}}(t)}
\end{equation}

\subsection*{5.3 Skala Tingkat Aktivitas Makan \& Matriks Keputusan}
\begin{table}[h!]
\centering
\small
\setlength{\tabcolsep}{3pt}
\caption{Matriks Klasifikasi Feeding Activity Index ($FAI$) dan Keputusan Kendali Feeder}
\label{tab:fai-matrix}
\begin{tabular}{|c|p{3.2cm}|p{5.5cm}|p{5.2cm}|}
\hline
\textbf{Rentang Skor $FAI$} & \textbf{Status Perilaku} & \textbf{Kondisi Visual Kolam} & \textbf{Aksi Sistem Kendali} \\
\hline
$FAI < 2,0$ & \textbf{Tenang / Kenyang} (\textit{Quiescent}) & Hanya semburan buih aerator normal. Permukaan air stabil. Bibit di dasar. & \textbf{Aktuator OFF}. Pakan tidak boleh ditebar untuk mencegah pembusukan pelet. \\
\hline
$2,0 \le FAI < 4,0$ & \textbf{Aktivitas Moderat} (\textit{Searching}) & Beberapa ekor bibit berenang naik, muncul riak mikro sporadis kecil. & \textbf{Standby Monitoring}. Perekaman intensif, belum memicu aktuator. \\
\hline
$FAI \ge 4,0$ & \textbf{Sangat Lapar / Agresif} (\textit{Frenzy}) & Benih serentak memecah permukaan air, riak dan semburan buih meluas masif. & \textbf{TRIGGER SERVO ON!} Dispensi 1 sub-porsi (50--100 g) jika kuota tersedia. \\
\hline
\end{tabular}
\end{table}

\newpage

\section*{6. Desain Sistem Terintegrasi \& Logika Kendali Aktuator Servo (PCA9685)}

\alertbox{LOGIKA FINITE STATE MACHINE (FSM) FEEDER ADAPTIF TERIKAT KUOTA HARIAN}{
\begin{itemize}
    \item \textbf{State 1 (Monitoring Pasif)}: Kamera merekam 10 FPS 720p. Nilai $FAI_t$ dievaluasi secara kontinu. Baseline aerator ($\mu_{\text{base}}, \sigma_{\text{base}}$) diperbarui terus-menerus.
    \item \textbf{State 2 (Verifikasi Kuota Harian)}: Ketika $FAI_t \ge 4,0$ selama $\ge 0,5\text{ detik}$, sistem mengecek sisa kuota:
    $$\text{Sisa Kuota} = \text{Kuota Harian Petani} - \text{Total Pakan Tersalurkan Hari Ini}$$
    Jika $\text{Sisa Kuota} \le 0\text{ gram}$, aktuator dikunci rapat dan pakan dibatalkan untuk mencegah keracunan amonia.\\[2pt]
    Jika $\text{Sisa Kuota} > 0\text{ gram}$, sistem melanjutkan ke \textbf{State 3}.
    \item \textbf{State 3 (Penggerak Servo PCA9685)}: Raspberry Pi mengirim sinyal I2C ke modul PCA9685 untuk memutar servo dari $0^\circ$ ke $90^\circ$ selama $2,0\text{ detik}$ (menjatuhkan $50\text{ gram}$ pakan pelet benih), lalu menutup kembali ke $0^\circ$. Catatan akumulasi pakan ditambah $+50\text{ gram}$.
    \item \textbf{State 4 (Jeda Konsumsi / Cooldown)}: Aktuator servo dikunci total selama \textbf{3 hingga 5 menit}. Jeda ini memberi waktu bagi benih lele untuk menghabiskan pakan yang ada dan mencegah cipratan makan memicu dispensi berulang kali.
\end{itemize}
}

\subsection*{6.1 Rangkaian Hardware Aktuator Servo PCA9685}
\begin{itemize}
    \item \textbf{Antarmuka I2C}: Pin SDA (GPIO 2) dan SCL (GPIO 3) Raspberry Pi 5 terhubung ke driver PCA9685 (alamat I2C \texttt{0x40}).
    \item \textbf{Isolasi Catu Daya}: Modul logika PCA9685 ditenagai 3.3V dari Raspberry Pi, namun rel daya motor servo ($V+$) disuplai adaptor eksternal $5\text{V } 3\text{A}$ terpisah guna mencegah \textit{voltage sag} yang dapat mereset Raspberry Pi saat servo bergerak.
    \item \textbf{Sinyal PWM ($50\text{ Hz}$)}: Tutup ($0^\circ$) = pulsa $1,0\text{ ms}$; Buka ($90^\circ$) = pulsa $2,0\text{ ms}$.
\end{itemize}

\section*{7. Rencana Pengujian Lapangan 1 Bulan \& Metrik Evaluasi Akademik}

Pengujian lapangan membandingkan dua kolam pendederan lele selama \textbf{30 hari kalender}:
\begin{itemize}
    \item \textbf{Kolam Perlakuan (Sistem Adaptif Skripsi)}: Pakan ditebar otomatis berbasis deteksi riak $FAI$ Raspberry Pi dan servo PCA9685 dengan batas kuota harian.
    \item \textbf{Kolam Kontrol (Metode Konvensional Petani)}: Pakan ditebar manual 3 kali sehari (07.00, 13.00, 19.00) dengan porsi tetap tanpa observasi riak.
\end{itemize}

\subsection*{7.1 Parameter Kinerja Biologis \& Budidaya}
\begin{enumerate}
    \item \textbf{Laju Pertumbuhan Bobot Harian (\textit{Average Daily Gain / ADG})}:
    \begin{equation}
    ADG = \frac{W_t - W_0}{t} \quad (\text{gram/hari})
    \end{equation}
    \item \textbf{Rasio Konversi Pakan (\textit{Feed Conversion Ratio / FCR})}:
    \begin{equation}
    FCR = \frac{\text{Total Bobot Pakan Diberikan (g)}}{\text{Pertambahan Total Biomassa Ikan (g)}}
    \end{equation}
    \textit{Target}: Kolam perlakuan menghasilkan FCR lebih efisien ($1,0\text{--}1,1$) dibanding kontrol ($1,3\text{--}1,5$).
    \item \textbf{Tingkat Kelangsungan Hidup (\textit{Survival Rate / SR})}:
    \begin{equation}
    SR = \frac{N_t}{N_0} \times 100\% \quad (\text{Target } \ge 85\%)
    \end{equation}
    \item \textbf{Efisiensi Penghematan Pakan (\textit{Feed Saving Efficiency})}:
    \begin{equation}
    \text{Efisiensi} = \frac{\text{Konsumsi Pakan Kontrol} - \text{Konsumsi Pakan Adaptif}}{\text{Konsumsi Pakan Kontrol}} \times 100\%
    \end{equation}
\end{enumerate}

\subsection*{7.2 Parameter Kinerja Visi Komputer Edge Device}
\begin{itemize}
    \item \textbf{Akurasi Deteksi Lapar}: Precision, Recall, Accuracy, dan F1-Score terhadap rekaman video teranotasi manual (*ground truth*). Target akurasi $\ge 85\%$.
    \item \textbf{Tingkat Alarm Palsu (*False Positive Rate / FPR*)}: Target $< 5\%$ terhadap letupan buih aerator normal atau pantulan sinar matahari.
    \item \textbf{Latensi Pemrosesan Edge}: Waktu eksekusi satu frame pada CPU Raspberry Pi 5 ($< 70\text{ ms}$ untuk memastikan pipeline 10 FPS tanpa frame drop).
\end{itemize}

\section*{8. Kesimpulan \& Panduan Implementasi Praktik Mahasiswa}

\begin{enumerate}
    \item \textbf{Kelengkapan Ilmiah}: Dokumen ini telah merumuskan secara tuntas seluruh jawaban teknis dan biologis, siap dikutip dan diadaptasi ke dalam naskah proposal/skripsi Bab 2 dan Bab 3.
    \item \textbf{Kecukupan Literatur}: Peneliti \textbf{TIDAK PERLU lagi mencari artikel via Consensus}. Literatur 20 artikel CV dan 14 artikel pakan sudah lebih dari cukup dan komprehensif.
    \item \textbf{Langkah Konkret Selanjutnya}:
    \begin{itemize}
        \item Cetak dokumen PDF terlampir sebagai bahan belajar intensif dan diskusi dengan dosen pembimbing.
        \item Masukkan polygon ROI masking pada kode kamera Raspberry Pi untuk membuang dinding terpal biru kolam.
        \item Lakukan kalibrasi nilai baseline aerator ($\mu_{\text{base}}, \sigma_{\text{base}}$) pada video riil selama 10 menit.
        \item Sambungkan modul driver PCA9685 ke Raspberry Pi 5 untuk menguji pergerakan sudut servo katup pakan.
    \end{itemize}
\end{enumerate}

\vspace{0.2cm}
\rule{\textwidth}{0.6pt}
\begin{center}
    \footnotesize
    Dokumen disiapkan khusus untuk Khoirotun Nisa' (NIM: 2341720057) --- D4 Teknik Informatika, Jurusan TI, Politeknik Negeri Malang.\\
    Dosen Pembimbing: Dr. Ulla Delfana Rosiani, ST., MT. $\bullet$ Tahun Akademik 2026/2027.
\end{center}

\end{document}
